% =========================================================
% Field and Order Axioms
% =========================================================

\section{Field and Order Axioms}
\label{sec:field-order-axioms}

\subsection{Definitions and Theorems}

The rational numbers are not merely a collection of equivalence classes.  Once addition and
multiplication are shown to be well-defined, $\mathbb{Q}$ becomes a field.  In addition, it carries a
natural order relation compatible with the field operations.  Together these two structures make
$\mathbb{Q}$ an \emph{ordered field}.

\paragraph{Field axioms.}
A field is a set $F$ with operations $+$ and $\cdot$ satisfying closure, associativity,
commutativity, additive and multiplicative identities, additive inverses, multiplicative inverses for
nonzero elements, and the distributive law.

\begin{remark}
The rational numbers satisfy these axioms under the usual operations inherited from their
construction.
\end{remark}

\paragraph{Order axioms.}
An ordered field is a field together with a relation $<$ satisfying trichotomy, transitivity,
compatibility with addition, and compatibility with multiplication by positive elements.

\begin{remark}
For many proofs it is more efficient to define order using the set of positive elements rather than to
take $<$ as primitive.
\end{remark}

\begin{definition}[Positive cone]
\label{def:positive-cone}
Let $F$ be a field.  A subset $P\subset F$ is called a \emph{positive cone} if it satisfies:
\begin{enumerate}
  \item if $a,b\in P$, then $a+b\in P$;
  \item if $a,b\in P$, then $ab\in P$;
  \item for every $a\in F$, exactly one of the following holds:
  \[
  a\in P,
  \qquad
  a=0,
  \qquad
  -a\in P.
  \]
\end{enumerate}
Define the order relation by
\[
a<b \iff b-a\in P.
\]
\end{definition}

\begin{remark}
This viewpoint makes the interaction between algebra and order especially transparent.  Statements
about inequalities can be translated into statements about positive elements and then handled
algebraically.
\end{remark}

\subsection{Immediate Consequences of the Positive Cone Axioms}

\begin{theorem}
\label{thm:squares-nonnegative}
For every $a\in F$,
\[
a^2\ge 0.
\]
\end{theorem}

\begin{proof}
By trichotomy, either $a\in P$, or $a=0$, or $-a\in P$.  In the first case $a^2\in P$ because
$P$ is closed under multiplication.  In the second case $a^2=0$.  In the third case
\[
(-a)(-a)=a^2\in P.
\]
Thus $a^2$ is positive or zero.
\end{proof}

\begin{theorem}
\label{thm:inverse-positive}
If $a>0$, then $a^{-1}>0$.
\end{theorem}

\begin{proof}
Suppose $a>0$.  If $a^{-1}<0$, then $-a^{-1}>0$.  Multiplying two positive elements gives
\[
a(-a^{-1})=-1>0.
\]
If also $1>0$, then trichotomy fails; so the assumption $a^{-1}<0$ is impossible.  Therefore
$a^{-1}>0$.
\end{proof}

\begin{theorem}
\label{thm:multiply-positive}
If $a<b$ and $c>0$, then
\[
ac<bc.
\]
\end{theorem}

\begin{proof}
From $a<b$ we have $b-a>0$.  Multiplying by $c>0$ gives
\[
c(b-a)>0.
\]
Thus
\[
bc-ac>0,
\]
which is equivalent to $ac<bc$.
\end{proof}

\begin{theorem}
\label{thm:multiply-negative}
If $a<b$ and $c<0$, then
\[
ac>bc.
\]
\end{theorem}

\begin{proof}
If $c<0$, then $-c>0$.  By the previous theorem,
\[
a(-c)<b(-c).
\]
This is
\[
-ac<-bc.
\]
Multiplying both sides by $-1$ reverses the inequality and yields $ac>bc$.
\end{proof}

\subsection{Positivity of $1$ and of the Natural Numbers}

\begin{theorem}
\label{thm:one-positive}
In an ordered field,
\[
1>0.
\]
\end{theorem}

\begin{proof}
By trichotomy exactly one of the following holds:
\[
1>0,
\qquad
1=0,
\qquad
-1>0.
\]
The second case is impossible in a field.  If $-1>0$, then multiplying two positive elements gives
\[
(-1)(-1)=1>0.
\]
Thus both $1$ and $-1$ would be positive, contradicting trichotomy.  Hence $1>0$.
\end{proof}

\begin{corollary}
\label{cor:n-positive}
For every natural number $n$,
\[
n>0.
\]
\end{corollary}

\begin{proof}
Each natural number is a finite sum of copies of $1$.  Since $1>0$ and the positive cone is closed
under addition, every such sum is positive.
\end{proof}

\begin{remark}[Archimedean property versus completeness]
\label{rem:arch-v-complete}
The positivity of $1$ anchors the order structure and leads naturally to the Archimedean property.
The Archimedean property holds in both $\mathbb{Q}$ and $\mathbb{R}$.  Completeness does not.  That is
why the rational numbers, although already an ordered field, still fail to capture all limits.
\end{remark}

% ---------------------------------------------------------
% Conceptual significance
% ---------------------------------------------------------

\begin{remark}[Conceptual significance]

The positivity of $1$ anchors the order structure of the field.  It
ensures that the natural numbers inherit a consistent ordering within
the field.

This connection between the algebraic structure of the field and the
order structure is essential for developing analysis, particularly in
the study of limits, bounds, and completeness.

\end{remark}


\begin{remark}

The Archimedean property holds in both the rational numbers
$\mathbb{Q}$ and the real numbers $\mathbb{R}$.

However, the rational numbers are not complete.  There exist
Cauchy sequences of rational numbers that do not converge to a
rational number.

Thus the real numbers can be characterized as the
\emph{complete Archimedean ordered field}.

\end{remark}



\begin{remark}[Archimedean property vs.\ completeness]

The Archimedean property holds in both the rational numbers
$\mathbb{Q}$ and the real numbers $\mathbb{R}$.

However, the rational numbers are \emph{not complete}.  There exist
Cauchy sequences of rational numbers that do not converge to a
rational limit.  For example, the sequence of rational approximations

\[
1,\;1.4,\;1.41,\;1.414,\;1.4142,\dots
\]

is Cauchy in $\mathbb{Q}$ but converges to $\sqrt{2}$, which is not
rational.

The real numbers $\mathbb{R}$ are obtained by adjoining these missing
limits.  Consequently, $\mathbb{R}$ satisfies the completeness
property: every Cauchy sequence converges.

A fundamental theorem of analysis states that the real numbers are
the \emph{unique complete Archimedean ordered field}.

\end{remark}



% =========================================================
% Mediant Property vs Archimedean Property
% =========================================================

\subsection{Mediant Property vs.\ Archimedean Property}

% ---------------------------------------------------------
% Teacher exposition
% ---------------------------------------------------------

\begin{remark}[Teacher exposition]
\label{rem:mediant-archimedean-teacher}

Two important structural properties of the rational numbers are the
\emph{mediant property} and the \emph{Archimedean property}.  Although
both describe fundamental features of $\mathbb{Q}$, they capture very
different aspects of the number system.

The mediant property concerns the refinement of rational intervals:
given two rational numbers, it provides a constructive way to produce
another rational number lying between them.

The Archimedean property, in contrast, describes the large-scale
behavior of the number system.  It states that the natural numbers
grow beyond any fixed element of the field.

Thus the mediant property describes the \emph{local structure} of
rational numbers, while the Archimedean property describes their
\emph{global growth behavior}.

\end{remark}

% ---------------------------------------------------------
% Mediant property
% ---------------------------------------------------------

\begin{definition}[Mediant]
\label{def:mediant}

Let

\[
\frac{a}{b}, \frac{c}{d} \in \mathbb{Q},
\qquad b,d>0.
\]

The \emph{mediant} of these fractions is

\[
\frac{a+c}{b+d}.
\]

\end{definition}

\begin{theorem}[Mediant Inequality]
\label{thm:mediant-inequality}

If

\[
\frac{a}{b} < \frac{c}{d},
\]

then

\[
\frac{a}{b}
<
\frac{a+c}{b+d}
<
\frac{c}{d}.
\]

\end{theorem}

\begin{remark}

The mediant provides a constructive method for inserting a new
rational number between two existing rational numbers.  Repeated
application of the mediant operation produces infinitely many
rational numbers inside any rational interval.

Thus the mediant property illustrates the \emph{density} of the
rational numbers.

\end{remark}

% ---------------------------------------------------------
% Archimedean property
% ---------------------------------------------------------

\begin{definition}[Archimedean Property]
\label{def:archimedean}

An ordered field $F$ is said to satisfy the \emph{Archimedean property}
if for every element $x \in F$ there exists a natural number $n$ such
that

\[
n > x.
\]

\end{definition}

\begin{remark}

An equivalent formulation is the following: for every positive number
$x$ there exists a natural number $n$ such that

\[
\frac{1}{n} < x.
\]

This version emphasizes that the reciprocals of natural numbers can
be made arbitrarily small.

\end{remark}

\begin{remark}

The rational numbers $\mathbb{Q}$ satisfy the Archimedean property.
The real numbers $\mathbb{R}$ also satisfy this property.

The Archimedean property rules out the existence of infinitely large
numbers or infinitesimal elements within the field.

\end{remark}

% ---------------------------------------------------------
% Conceptual comparison
% ---------------------------------------------------------

\begin{remark}[Conceptual comparison]

The mediant property and the Archimedean property describe two
different structural features of the rational numbers.

\begin{center}
\begin{tabular}{lll}
\textbf{Property} & \textbf{Description} & \textbf{Scope} \\
\toprule
Mediant property & Inserts new rationals between two rationals & Local \\
Archimedean property & Natural numbers exceed every element & Global
\end{tabular}
\end{center}

The mediant property concerns the refinement of rational intervals,
while the Archimedean property concerns the overall scale of the
number system.

\end{remark}

% ---------------------------------------------------------
% Final structural remark
% ---------------------------------------------------------

\begin{remark}

The rational numbers possess both the mediant property and the
Archimedean property.  Nevertheless, they still contain gaps.  For
example,

\[
1.4 < 1.41 < 1.414 < \sqrt{2} < 1.415 < 1.42.
\]

Although rational numbers can approximate $\sqrt{2}$ arbitrarily
closely, no rational number equals $\sqrt{2}$.

This shows that $\mathbb{Q}$ is dense and Archimedean but not
complete.  The real numbers $\mathbb{R}$ are obtained by adjoining
the missing limits, producing a \emph{complete Archimedean ordered
field}.

\end{remark}

\begin{theorem}
\[
1>0
\]
\end{theorem}

\begin{proof}
By trichotomy exactly one of the following holds:
\[
1>0, \qquad 1=0, \qquad -1>0.
\]

The second is impossible in a field.  
If $-1>0$ then multiplying two positive numbers gives

\[
(-1)(-1)=1>0.
\]

But then both $1$ and $-1$ would be positive, contradicting the
trichotomy property.
\end{proof}

\begin{corollary}
For every $n\in\mathbb{N}$,
\[
n>0.
\]
\end{corollary}

\begin{proof}
Since $1>0$ and $P+P\subseteq P$, repeated addition yields
\[
1+1+\cdots+1=n>0.
\]
\end{proof}



















